% Source: ACM_Templates.pdf, PDF pp. 29-31.
% Corrections and validation: docs/validation.md.
% Additional reference: https://cp-algorithms.com/combinatorics/binomial-coefficients.html
\section{Lucas 与扩展 Lucas}
\label{algo:lucas}
\subsection*{Lucas 定理：素数模数}
当 $p$ 为素数，将 $n,k$ 写成 $p$ 进制：$n=\sum n_ip^i$、$k=\sum k_ip^i$，则
\[
    \binom nk\equiv\prod_i\binom{n_i}{k_i}\pmod p.
\]
由 $(1+x)^p\equiv1+x^p\pmod p$，可把 $(1+x)^n$ 按进制拆开，比较 $x^k$ 的系数得到该式。某一位 $k_i>n_i$ 时，该项为 $0$，整个结果也为 $0$。

预处理 $0,\ldots,p-1$ 的阶乘与阶乘逆元，时间 $O(p+\log p)$、空间 $O(p)$，每次查询 $O(\log_p(n+1))$。关键是预处理到 $p-1$；不是按原始的大 $n$ 建表，也不能把 $p!\equiv0$ 当成可逆数。本模板面向较小的素数 $p$，例如 $p\leq10^6$；这属于资源约束，定理本身不受此上限限制。

\subsection*{使用教程}
先复制快速幂与 \InlineCode{Comb}。\InlineCode{Lucas::init(p)} 只针对模数预处理，之后用 \InlineCode{C(n,k)} 查询非负 64 位参数。
\begin{lstlisting}[style=mtycpp,numbers=none]
Lucas lc;
lc.init(101);
long long ans=lc.C(1000000000000LL,12345);
\end{lstlisting}

\needspace{7\baselineskip}
\subsection*{模板代码}
\CodeFile{code/number-theory/lucas.cpp}

\subsection*{扩展 Lucas：任意小模数}
把 $m$ 分解为两两互质的素数幂 $m=\prod_i q_i$，$q_i=p_i^{\alpha_i}$。先分别求模 $q_i$ 的组合数，再用中国剩余定理合并。因为阶乘可能含有 $p_i$，必须把“$p_i$ 的指数”和“去掉这些因子后可逆的部分”分开。

固定 $q=p^\alpha$，令 $e(n)=v_p(n!)$、$F(n)=n!/p^{e(n)}\bmod q$，则
\[
    \binom nk\equiv
    F(n)F(k)^{-1}F(n-k)^{-1}
    p^{e(n)-e(k)-e(n-k)}\pmod q.
\]
若指数至少为 $\alpha$，答案直接为 $0$。其余三个 $F$ 都不含因子 $p$，因而可以对 $q$ 求逆。

设 $U(t)=\prod_{1\leq j\leq t,\,p\nmid j}j\bmod q$，预处理 $U(0),\ldots,U(q)$。把非 $p$ 倍数按长度 $q$ 分块，把 $p$ 的倍数除去一个 $p$ 后递归，得到
\[
    F(n)=U(q)^{\lfloor n/q\rfloor}U(n\bmod q)F(\lfloor n/p\rfloor)\pmod q.
\]
这也是代码中循环求去因子阶乘的依据。

本实现预处理 $O(\sqrt m+\sum_i q_i)$、空间 $O(\sum_i q_i)$。每层还用快速幂计算完整块的贡献，因此保守查询上界为
\[
    O\!\left(\sum_i\log_{p_i}(n+1)\log(n+m)+s\log m\right),
\]
其中 $s$ 为不同素因子数；不能把这份实现写成每层常数时间。模数虽可为合数，仍需足够小以存储各素数幂的表，建议 $m\leq2\times10^6$，并结合内存限制调整。

\subsection*{使用教程}
先复制快速幂、扩欧、逆元及第~\ref{algo:excrt}~节的 \InlineCode{ExCRT}。\InlineCode{ExLucas::init(m)} 对正模数预处理一次；\InlineCode{C(n,k)} 可反复查询，$k$ 非法或 $m=1$ 时返回 $0$。
\begin{lstlisting}[style=mtycpp,numbers=none]
ExLucas ex;
ex.init(72);
long long ans=ex.C(1000000,123);
\end{lstlisting}

\needspace{7\baselineskip}
\subsection*{模板代码}
\CodeFile{code/number-theory/exlucas.cpp}
